% Related work grouped by energy accounting, visual transmission, and
% adaptive inference. Comparisons identify decision variables and decoders.
\section{Related Work}\label{sec:related}

\subsection{Energy-Efficient Semantic Communication}
Energy-efficient semantic communication considers the cost of extracting
and processing information as well as transmitting it.
Zheng \etal~\cite{zheng2024aerialee} study aerial-aided edge networks
and develop incentive and distributed learning mechanisms for semantic
transmission and coder updates. PSCom~\cite{dai2025pscom} minimizes
communication and computation energy under constraints on latency,
power, bandwidth, and semantic compression. Rate-splitting designs
jointly allocate communication resources and semantic compression
levels~\cite{xu2025rsmagreen}. Related formulations address UAV networks,
including flight energy~\cite{wang2026rsma}, and visible-light systems
with knowledge-update overhead~\cite{zhao2026vlc}. These studies show
why savings in transmitted data must be evaluated together with the
computational cost of semantic processing.

For VLM-based services, the inference workload also affects the energy
account. Li \etal~\cite{li2026onboard} consider onboard VQA and optimize
image resolution, transmission power, and UAV trajectory under accuracy
constraints. Device-level measurements by Zhan
\etal~\cite{zhan2026seeing} identify autoregressive output generation
as a major energy cost in the VLM configurations they evaluate. Their
results indicate that reducing visual input alone may provide limited
inference-energy savings. These findings motivate accounting for the
processing required after transmission. Our study evaluates a
representation choice that also changes the answering computation:
symbolic decoding for detector information or VLM inference for an image.
UAV-side and edge-side costs are distinguished because they draw on
different energy supplies.

\subsection{Task-Oriented Visual Communication}
Learned image transmission provides one way to reduce communication
cost. Deep joint source--channel coding (DJSCC) maps images directly to
channel symbols and reconstructs them at the
receiver~\cite{bourtsoulatze2019djscc}. Attention-based coding adapts
to channel conditions~\cite{yang2023adjscc}, while predictive adaptive
coding uses estimated reconstruction quality to select the transmission
rate~\cite{zhang2023padc}. These methods improve the delivery of an
image representation. Task-oriented schemes additionally use the
requirements of the downstream task to determine what should be sent.

For wireless VQA, MU-DeepSC learns correlated image and question
representations with a jointly designed
transceiver~\cite{xie2022mudeepscwcl,xie2021mudeepsc}.
Goal-oriented VQA communication ranks bounding boxes and scene-graph
triplets using the question and employs an answer
reasoner~\cite{liu2024gosg}. Its comparison of representations across
question categories is directly relevant to evidence selection.
VideoQA-SC uses bandwidth-adaptive coding to answer video questions
from transmitted semantics without reconstructing the
video~\cite{guo2025videoqasc}. These approaches already evaluate
communication through the quality of the resulting answers, rather
than requiring accurate pixel reconstruction.

Other work adapts the representation or its protection to the task.
Park and Yoon~\cite{park2025transmit} compare object information,
layouts, and scene graphs across visual tasks and filter redundant
graph content. TONIC~\cite{liu2026tonic} assigns unequal protection
according to token relevance and combines receiver-side gating with
token completion for image classification. With pretrained multimodal
models, query-oriented coding selects features for query-relevant image
reconstruction~\cite{huang2026query}, and Du
\etal~\cite{jiang2025lmmvn} prioritize image regions and patch
transmission for a LLaVA-based vehicle assistant. These works motivate
task-dependent representation and transmission choices. The specific
choice examined here is between two complete answering paths within
a structured VQA workload. Each path has its own representation,
decoder, accuracy, and energy cost. Consequently, a smaller payload
may save both transmission energy and VLM inference, depending on
which decoder answers the question.

\subsection{Adaptive Inference and Evidence Selection}
Adaptive inference changes where computation takes place or how much
information reaches the server. Neurosurgeon~\cite{kang2017neurosurgeon}
partitions neural-network execution between a mobile device and the
cloud to reduce latency and mobile energy consumption. Information-bottleneck
methods learn compact representations for task-oriented edge
inference~\cite{shao2022taskoriented}. For VQA, INAR-VL
uses image and text complexity signals to route requests between
complementary edge and cloud VLMs~\cite{sabanovic2026inarvl}.
This provides an input-dependent choice of model and execution location.

Adaptive evidence delivery offers a related control variable.
Song and Kim~\cite{song2026collaborative} first send a thumbnail for
server-side VLM inference. The server uses output uncertainty to decide
whether to request an additional image region. Progressive semantic
communication compresses visual tokens into refinable representations
that can be supplied to off-the-shelf VLMs at different communication
budgets~\cite{hsu2026progressive}. SAGE~\cite{choi2026sage} selects
complementary visual evidence under a hard uplink budget and evaluates
the resulting classification accuracy. In remote sensing VQA,
Prompt-RSVQA converts extracted visual context into text for a language
model~\cite{chappuis2022promptrsvqa}. Thus, adaptive evidence selection
and structured context are established ways to support inference.

Our design combines sender-side evidence selection with a choice of
answering mechanism. It selects detector tokens for a symbolic decoder
or an image for a frozen VLM before transmitting the selected payload.
The decision does not require an initial receiver VLM pass or a request
for additional evidence. We assess its energy--accuracy trade-off using
logged-outcome replay, GPU power measurements, and models of transmission
and onboard detection costs. This distinguishes evidence-level routing
from model-placement decisions and from adapting the input to a VLM
that still runs on every query.
